\scanpage{104}
\mspage{\textcircled{7}}
i.e. on a une équivalence de Cat. 2-cofibrées
\[
T_{\deleted{\unsure{\lambda}}}^{o}\simeq(T_{\deleted{\unsure{\lambda}}})_{\unsure{o|\circ}}^{\unsure{o}}\quad\text{(passage à l'opposé)}\tag*{(7.1)}
\]
On a $(7.\unsure{2|3})$ \unsure{!}

On trouve ainsi une 2-équivalence
\[
(\lambda\text{-types})\simeq(\lambda^o\text{-types})\tag*{(7.3)}
\]
(et kif kif sur $I$
\[
(\lambda\text{-}I\text{-types})\simeq(\lambda^o\text{-}I\text{-types}))\tag*{(7.4.)}
\]
compatible, via les cat. représentatives, à
\[
\begin{aligned}
(\mathrm{Cat}_\lambda)&\simeq(\mathrm{Cat}_{\lambda^o})\\
C&\longmapsto C^o
\end{aligned}\tag*{(7.6)}
\]

\fbox{8} Cas où $\underset{\rightarrow}{\lambda}=\varnothing$ (structures définies via $\varprojlim$)

\inserted{$T$ désignera une $\lambda$-structure.}
8.1. Pour tout $C\in\mathrm{Ob}\,\mathrm{Cat}_{\lambda}$, les \deleted{limites}\inserted{$\varprojlim$} de type $\lambda$ existent dans $T(C)$.
plus précisément, \deleted{\unsure{les lim}} \inserted{tout type $d$} de $\varprojlim$ qui existent dans $C$
existent aussi dans $T(C)$. [Ecrire $T(C)\simeq\underline{\mathrm{Hom}}_\lambda(R_T,C)$,
et noter que \deleted{\unsure{toutes les}} l'hyp. sur $C$ implique que $\underline{\mathrm{Hom}}(R_T,C)$
\deleted{est stable} a des $\varprojlim$ de type $d$ \inserted{se calculant argt par $\underline{\text{argument}}$}, et que \deleted{\unsure{les f}} $\underline{\mathrm{Hom}}_\lambda(R_T,C)$
est \deleted{$\underline{\mathrm{Hom}}_\lambda$} est stable par lesdites $\varprojlim$. De plus, pour tt
$\unsure{r|z|\xi}\in R_T$, \deleted{\unsure{Ev}} $r_C:T(C)\longrightarrow C$ commute aux $\varprojlim$ de type $d$

8.1.1. Plus généralement ($\lambda$ \uline{quelconque}, sans néc. $\underset{\rightarrow}{\lambda}=\varnothing$)
Si \inserted{$C\in\mathrm{Ob}\,\mathrm{Cat}\,C_\lambda$ \unsure{$d$} et si} $d$ \deleted{\unsure{soit}} est un type de diagrammes tel que

a) $C$ admet des $\varprojlim$ de type $d$ (p.ex. $d\in\underset{\leftarrow}{\lambda}$)

b) dans $C$, les $\varprojlim$ de type $d$ commutent aux
$\varinjlim$ de type $\in\underset{\rightarrow}{\lambda}$

Alors \deleted{\unsure{toutes}} $T(C)$ admet des $\varprojlim$ de type $d$, et
\deleted{\unsure{si} \illegible{1} $R_T^\lambda$ alors \unsure{se}} écrivant $T(C)\simeq\underline{\mathrm{Hom}}_\lambda(R_T^\lambda,C)$,
ces limites se calculent ``argument par argument''.

% slanted
\margin{de plus si $C'$ est comme $C$ et si $C\longrightarrow C'$ commute \deleted{$T(C)\to T(C')$} aux $\varprojlim$ de type $d$, alors $T(C)\to T(C')$ \deleted{\unsure{est}}. Kif-Kif pour $T\to T'$, \unsure{alors} $T(C)\to T'(C)$ commute aux $\varprojlim$ de type \unsure{$d$}.}

toute \deleted{\unsure{flè}} $\lambda$-construction $\rho$ rel. à $T$, le foncteur
correspondant
\[
\rho:T(C)\longrightarrow C
\]
% archive sheet number: 103
